\section{地震・弾性波探査}

\subsection{平成23年度 問29}
\label{q:h23-reg-applied-29}

\noindent\textbf{原典：} sources/H23/17\_applied\_science.pdf、PDF物理10頁。

\textbf{公式正答：} 5（single）。

\paragraph{問題の要旨} 物理探査の観測対象と装置を比較し、海上音波探査の受振器を誤って説明した記述を選ぶ。

\Needspace{4\baselineskip}\paragraph{確認した条件・記述（独自要約）}\mbox{}\par

\noindent\textbf{選択方向：} 誤り又は最も不適切・関係の薄いものを1つ\par

\paragraph{問われる核心} 観測量と受振器の対応

\paragraph{必要知識} ハイドロフォンの圧力応答、地上速度計、相対重力計、MT、電磁波の反射。

\paragraph{考え方・正答を選ぶ理由（査読待ち初稿）} 測定量を対照する。重力は加速度、MTは自然電磁応答、レーダは電磁反射である。水中の流体はせん断波を伝えず、海上の受振系では音圧を検出するハイドロフォンを用いる。速度計を防水加工するという第5肢の説明が誤りである。

\Needspace{4\baselineskip}\paragraph{選択肢の検討（査読待ち初稿）}\mbox{}\par

\noindent\textbf{1：適切} 磁気から鉄類の有無と概略位置を探るが、細かな形状や製品種類は判別できない。通常の磁気異常だけでは磁化・形状・深度が非一意であり、細かな製品識別までは確定できないという探査上の限界を述べる。現代の補助情報を用いる識別可能性全般を否定する命題とは解さない。\par

\noindent\textbf{2：適切} MTは自然の地磁気・地電流から深部比抵抗を探る。自然変動の水平電場・磁場の応答から電気抵抗率を推定する方法である。\par

\noindent\textbf{3：適切} 相対重力計は吊った錘のばね伸びを利用する。ばね定数と伸びの変化から重力の差を測る。絶対値を直接定める落下式とは区別する。\par

\noindent\textbf{4：適切} 地中レーダは概ね10 MHz以上の電磁波で極浅部を探る。原典が述べる通常の浅部レーダの周波数帯と原理に対応する。到達深度は周波数だけでなく導電率や減衰にも依存する。\par

\noindent\textbf{5：不適切} ストリーマケーブルには防水加工した速度計を多数入れる。水中音波を測る受振器は通常ハイドロフォンであり、圧力変動を検出する。地上用速度計の防水化と同一ではない。\par

\paragraph{再利用する考え方} 観測量と受振器の対応を、測定量・推定対象・基準面・近似条件に分けて理解する。名称だけで探査法を対応させない。

\paragraph{共通理論} \hyperref[theory:exploration:observables]{第\ref*{theory:exploration:observables}節}, \hyperref[theory:exploration:electromagnetic]{第\ref*{theory:exploration:electromagnetic}節}, \hyperref[theory:gravity:measurement]{第\ref*{theory:gravity:measurement}節}

\par 参考：\cite{usgs_marine_seismic,usgs_mt_metadata}。

\par 難易度 3（校準初稿）、優先度未評価。

\subsection{平成23年度 問32}
\label{q:h23-reg-applied-32}

\noindent\textbf{原典：} sources/H23/17\_applied\_science.pdf、PDF物理11頁。

\textbf{公式正答：} 3（single）。

\par\textit{原典本文の校正と独自解説は未作成。分類のみの登録項目。}

\subsection{平成25年度 問32}
\label{q:h25-reg-applied-32}

\noindent\textbf{原典：} sources/H25/17\_applied\_science.pdf、PDF物理12頁。

\textbf{公式正答：} 3（single）。

\par\textit{原典本文の校正と独自解説は未作成。分類のみの登録項目。}

\subsection{平成26年度 問20}
\label{q:h26-reg-applied-20}

\noindent\textbf{原典：} sources/H26/17\_applied\_science.pdf、PDF物理8頁。

\textbf{公式正答：} 1（single）。

\par\textit{原典本文の校正と独自解説は未作成。分類のみの登録項目。}

\subsection{平成29年度 問34}
\label{q:h29-reg-applied-34}

\noindent\textbf{原典：} sources/H29/17\_applied\_science.pdf、PDF物理12頁。

\textbf{公式正答：} 3（single）。

\par\textit{原典本文の校正と独自解説は未作成。分類のみの登録項目。}

\subsection{令和元年度再試験 問32}
\label{q:r01-ret-applied-32}

\noindent\textbf{原典：} sources/R01\_retest/17\_applied\_science.pdf、PDF物理11頁。

\textbf{公式正答：} 4（single）。

\par\textit{原典本文の校正と独自解説は未作成。分類のみの登録項目。}

\subsection{令和2年度 問32}
\label{q:r02-reg-applied-32}

\noindent\textbf{原典：} sources/R02/17\_applied\_science.pdf、PDF物理11頁。

\textbf{公式正答：} 2（single）。

\par\textit{原典本文の校正と独自解説は未作成。分類のみの登録項目。}

\subsection{令和3年度 問33}
\label{q:r03-reg-applied-33}

\noindent\textbf{原典：} sources/R03/17\_applied\_science.pdf、PDF物理12頁。

\textbf{公式正答：} 1（single）。

\par\textit{原典本文の校正と独自解説は未作成。分類のみの登録項目。}

\subsection{令和3年度 問34}
\label{q:r03-reg-applied-34}

\noindent\textbf{原典：} sources/R03/17\_applied\_science.pdf、PDF物理13頁。

\textbf{公式正答：} 2（single）。

\par\textit{原典本文の校正と独自解説は未作成。分類のみの登録項目。}

\subsection{令和5年度 問32}
\label{q:r05-reg-applied-32}

\noindent\textbf{原典：} sources/R05/17\_applied\_science.pdf、PDF物理11頁。

\textbf{公式正答：} 4（single）。

\par\textit{原典本文の校正と独自解説は未作成。分類のみの登録項目。}

\subsection{令和6年度 問32}
\label{q:r06-reg-applied-32}

\noindent\textbf{原典：} sources/R06/17\_applied\_science.pdf、PDF物理11頁。

\textbf{公式正答：} 1（single）。

\par\textit{原典本文の校正と独自解説は未作成。分類のみの登録項目。}

\section{重力・磁気探査}

\subsection{平成27年度 問22}
\label{q:h27-reg-applied-22}

\noindent\textbf{原典：} sources/H27/17\_applied\_science.pdf、PDF物理8頁。

\textbf{公式正答：} 2（single）。

\paragraph{問題の要旨} 地形、無限水平層、潮汐、ばねドリフト、緯度の重力効果から誤りを選ぶ。

\Needspace{4\baselineskip}\paragraph{確認した条件・記述（独自要約）}\mbox{}\par

\noindent\textbf{選択方向：} 誤り又は最も不適切・関係の薄いものを1つ\par

\paragraph{問われる核心} 無限平板の重力と重力補正の基準

\paragraph{必要知識} ニュートン引力積分、無限平板と有限岩体の違い、地形補正の符号、潮汐、ドリフト、正規重力。

\paragraph{考え方・正答を選ぶ理由（査読待ち初稿）} 無限水平層の上に観測点がある場合、各薄層の寄与を積分すると\(g_z=2\pi G\rho t\)である。層と測点の距離は消えるので、浅いほど大きいという第2肢が誤りである。第1肢は無限平板に対する地形差、第4肢は短時間の線形ドリフト近似として条件を明示して読む。

\Needspace{4\baselineskip}\paragraph{選択肢の検討（査読待ち初稿）}\mbox{}\par

\noindent\textbf{1：適切} 山・谷の地形効果はともに重力を小さくする。測点高まで無限平板があるというブーゲー基準に対する地形効果として読む。測点より上の山の引力は上向き、谷は基準平板の下向き引力を欠くため、いずれも下向き重力を減らす。任意の山と谷の絶対重力を比較する命題ではない。\par

\noindent\textbf{2：不適切} 密度と厚さ一定の無限水平層は浅いほど重力が大きい。無限平板の外部重力は2πGρtで深度に依存しない。有限の岩体では距離依存があるため、無限という条件が重要である。\par

\noindent\textbf{3：適切} 月と太陽の引力による潮汐補正が必要である。月・太陽による時間変動を地下構造による異常と区別する必要がある。\par

\noindent\textbf{4：適切} ばねの時間変化を一定と近似し閉合誤差を各点へ配分する。同一点への戻り測定の閉合差から時間に比例するドリフトを近似して補正する操作を述べる。一定の伸びは測定区間のドリフトを一定率とする通常の近似として解釈し、あらゆる時間変化が必ず線形とはしない。\par

\noindent\textbf{5：適切} 絶対重力と緯度から求めた正規重力との差をとるのが緯度補正である。緯度に依存する理論基準を差し引き、基準地球の変化を除く操作である。\par

\paragraph{再利用する考え方} 無限平板の重力と重力補正の基準を、測定量・推定対象・基準面・近似条件に分けて理解する。名称だけで探査法を対応させない。

\paragraph{共通理論} \hyperref[theory:gravity:slab]{第\ref*{theory:gravity:slab}節}, \hyperref[theory:gravity:corrections]{第\ref*{theory:gravity:corrections}節}, \hyperref[theory:gravity:measurement]{第\ref*{theory:gravity:measurement}節}

\par 参考：\cite{gsi_gravity_definitions}。

\paragraph{関連問題} \hyperref[q:r07-reg-applied-33]{R07-REG-APPLIED-33}

\par 難易度 3（校準初稿）、優先度未評価。

\subsection{平成29年度 問31}
\label{q:h29-reg-applied-31}

\noindent\textbf{原典：} sources/H29/17\_applied\_science.pdf、PDF物理10頁。

\textbf{公式正答：} 5（single）。

\paragraph{問題の要旨} 重力探査で用いる補正名のうち、主に反射法地震探査に属するものを選ぶ。

\Needspace{4\baselineskip}\paragraph{確認した条件・記述（独自要約）}\mbox{}\par

\noindent\textbf{選択方向：} 誤り又は最も不適切・関係の薄いものを1つ\par

\paragraph{問われる核心} 重力補正と地震走時補正の区別

\paragraph{必要知識} 運動・質量・潮汐・標高の重力効果、浅部層による走時差。

\paragraph{考え方・正答を選ぶ理由（査読待ち初稿）} 補正の対象量を並べると、第1〜4肢は重力加速度の変化を扱う。静補正は主に地震探査の走時をそろえるもので、重力探査の標準的補正名としては該当しない。

\Needspace{4\baselineskip}\paragraph{選択肢の検討（査読待ち初稿）}\mbox{}\par

\noindent\textbf{1：適切} エトベス補正。船や航空機の運動に伴う回転系の見掛加速度の差を補正する。\par

\noindent\textbf{2：適切} ブーゲー補正。基準面と測点の間の岩石質量の引力を除く。\par

\noindent\textbf{3：適切} 潮汐補正。月・太陽による時間変動を除く。\par

\noindent\textbf{4：適切} フリーエア補正。基準面からの高さに伴う重力変化を補正する。\par

\noindent\textbf{5：不適切} 静補正。浅部低速度層や標高による地震波の走時差を、共通の基準面へ補正する地震探査の用語である。\par

\paragraph{再利用する考え方} 重力補正と地震走時補正の区別を、測定量・推定対象・基準面・近似条件に分けて理解する。名称だけで探査法を対応させない。

\paragraph{共通理論} \hyperref[theory:gravity:corrections]{第\ref*{theory:gravity:corrections}節}, \hyperref[theory:exploration:observables]{第\ref*{theory:exploration:observables}節}

\par 参考：\cite{gsi_gravity_definitions}。

\par 難易度 3（校準初稿）、優先度未評価。

\subsection{平成30年度 問32}
\label{q:h30-reg-applied-32}

\noindent\textbf{原典：} sources/H30/17\_applied\_science.pdf、PDF物理11頁。

\textbf{公式正答：} 3（single）。

\paragraph{問題の要旨} 重力・磁気探査の基準面、磁化特性、移動観測補正と、別の観測分野の語を区別する。

\Needspace{4\baselineskip}\paragraph{確認した条件・記述（独自要約）}\mbox{}\par

\noindent\textbf{選択方向：} 誤り又は最も不適切・関係の薄いものを1つ\par

\paragraph{問われる核心} 重力・磁気探査の物理量と用語

\paragraph{必要知識} 等ポテンシャル面、残留磁化と誘導磁化、磁気秩序の温度依存、回転座標系。

\paragraph{考え方・正答を選ぶ理由（査読待ち初稿）} ジオイドとエトベス補正は重力、ケーニヒスベルガー比とキュリー温度は磁気に直接関係する。プランク関数は熱平衡放射のスペクトルを記述する関数であり、この候補群では重力・磁気探査の基本量に当たらない。したがってこの肢を選ぶ。

\Needspace{4\baselineskip}\paragraph{選択肢の検討（査読待ち初稿）}\mbox{}\par

\noindent\textbf{1：適切} ジオイド。基準となる重力等ポテンシャル面であり、重力場と標高の基準に関係する。\par

\noindent\textbf{2：適切} ケーニヒスベルガー比。自然残留磁化の大きさと現在の場による誘導磁化の大きさの比であり、磁気異常の解釈に関係する。\par

\noindent\textbf{3：不適切} プランク関数。熱平衡放射のスペクトルを記述する関数であり、この候補群では重力・磁気探査の基本量に当たらない。\par

\noindent\textbf{4：適切} キュリー温度。強磁性・フェリ磁性の自発磁化が消失する温度に対応し、磁性鉱物や磁気探査の解釈に関係する。\par

\noindent\textbf{5：適切} エトベス補正。移動観測系の運動に伴う見掛重力の差を補正する。\par

\paragraph{再利用する考え方} 重力・磁気探査の物理量と用語を、測定量・推定対象・基準面・近似条件に分けて理解する。名称だけで探査法を対応させない。

\paragraph{共通理論} \hyperref[theory:gravity:corrections]{第\ref*{theory:gravity:corrections}節}, \hyperref[theory:gravity:geoid]{第\ref*{theory:gravity:geoid}節}, \hyperref[theory:geomagnetism:magnetization]{第\ref*{theory:geomagnetism:magnetization}節}, \hyperref[theory:exploration:observables]{第\ref*{theory:exploration:observables}節}

\par 参考：\cite{gsi_gravity_definitions,kanaya_okuma2011}。

\paragraph{関連問題} \hyperref[q:r04-reg-applied-35]{R04-REG-APPLIED-35}

\par 難易度 3（校準初稿）、優先度未評価。

\subsection{令和4年度 問35}
\label{q:r04-reg-applied-35}

\noindent\textbf{原典：} sources/R04/17\_applied\_science.pdf、PDF物理12頁。

\textbf{公式正答：} 4（single）。

\paragraph{問題の要旨} 重力・磁気探査の基準面、磁化特性、移動観測補正と、別の観測分野の語を区別する。

\Needspace{4\baselineskip}\paragraph{確認した条件・記述（独自要約）}\mbox{}\par

\noindent\textbf{選択方向：} 誤り又は最も不適切・関係の薄いものを1つ\par

\paragraph{問われる核心} 重力・磁気探査の物理量と用語

\paragraph{必要知識} 等ポテンシャル面、残留磁化と誘導磁化、磁気秩序の温度依存、回転座標系。

\paragraph{考え方・正答を選ぶ理由（査読待ち初稿）} ジオイドとエトベス補正は重力、ケーニヒスベルガー比とキュリー温度は磁気に直接関係する。プランク関数は熱平衡放射のスペクトルを記述する関数であり、この候補群では重力・磁気探査の基本量に当たらない。したがってこの肢を選ぶ。

\Needspace{4\baselineskip}\paragraph{選択肢の検討（査読待ち初稿）}\mbox{}\par

\noindent\textbf{1：適切} エトベス補正。移動観測系の運動に伴う見掛重力の差を補正する。\par

\noindent\textbf{2：適切} ジオイド。基準となる重力等ポテンシャル面であり、重力場と標高の基準に関係する。\par

\noindent\textbf{3：適切} キュリー温度。強磁性・フェリ磁性の自発磁化が消失する温度に対応し、磁性鉱物や磁気探査の解釈に関係する。\par

\noindent\textbf{4：不適切} プランク関数。熱平衡放射のスペクトルを記述する関数であり、この候補群では重力・磁気探査の基本量に当たらない。\par

\noindent\textbf{5：適切} ケーニヒスベルガー比。自然残留磁化の大きさと現在の場による誘導磁化の大きさの比であり、磁気異常の解釈に関係する。\par

\paragraph{再利用する考え方} 重力・磁気探査の物理量と用語を、測定量・推定対象・基準面・近似条件に分けて理解する。名称だけで探査法を対応させない。

\paragraph{共通理論} \hyperref[theory:gravity:corrections]{第\ref*{theory:gravity:corrections}節}, \hyperref[theory:gravity:geoid]{第\ref*{theory:gravity:geoid}節}, \hyperref[theory:geomagnetism:magnetization]{第\ref*{theory:geomagnetism:magnetization}節}, \hyperref[theory:exploration:observables]{第\ref*{theory:exploration:observables}節}

\par 参考：\cite{gsi_gravity_definitions,kanaya_okuma2011}。

\paragraph{関連問題} \hyperref[q:h30-reg-applied-32]{H30-REG-APPLIED-32}, \hyperref[q:r06-reg-applied-34]{R06-REG-APPLIED-34}

\par 難易度 3（校準初稿）、優先度未評価。

\subsection{令和6年度 問34}
\label{q:r06-reg-applied-34}

\noindent\textbf{原典：} sources/R06/17\_applied\_science.pdf、PDF物理12頁。

\textbf{公式正答：} 3（single）。

\paragraph{問題の要旨} 重力・磁気探査の基準面、磁化特性、移動観測補正と、別の観測分野の語を区別する。

\Needspace{4\baselineskip}\paragraph{確認した条件・記述（独自要約）}\mbox{}\par

\noindent\textbf{選択方向：} 誤り又は最も不適切・関係の薄いものを1つ\par

\paragraph{問われる核心} 重力・磁気探査の物理量と用語

\paragraph{必要知識} 等ポテンシャル面、残留磁化と誘導磁化、磁気秩序の温度依存、回転座標系。

\paragraph{考え方・正答を選ぶ理由（査読待ち初稿）} ジオイドとエトベス補正は重力、ケーニヒスベルガー比とキュリー温度は磁気に直接関係する。AE（アコースティック・エミッション）は微小破壊等に伴う過渡的な弾性波の放出・計測である。岩石研究にも用いるが、重力や磁化の測定とは異なる物理量を扱う。したがってこの肢を選ぶ。

\Needspace{4\baselineskip}\paragraph{選択肢の検討（査読待ち初稿）}\mbox{}\par

\noindent\textbf{1：適切} ケーニヒスベルガー比。自然残留磁化の大きさと現在の場による誘導磁化の大きさの比であり、磁気異常の解釈に関係する。\par

\noindent\textbf{2：適切} ジオイド。基準となる重力等ポテンシャル面であり、重力場と標高の基準に関係する。\par

\noindent\textbf{3：不適切} AE（アコースティック・エミッション）。微小破壊等に伴う過渡的な弾性波の放出・計測である。岩石研究にも用いるが、重力や磁化の測定とは異なる物理量を扱う。\par

\noindent\textbf{4：適切} キュリー温度。強磁性・フェリ磁性の自発磁化が消失する温度に対応し、磁性鉱物や磁気探査の解釈に関係する。\par

\noindent\textbf{5：適切} エトベス補正。移動観測系の運動に伴う見掛重力の差を補正する。\par

\paragraph{再利用する考え方} 重力・磁気探査の物理量と用語を、測定量・推定対象・基準面・近似条件に分けて理解する。名称だけで探査法を対応させない。

\paragraph{共通理論} \hyperref[theory:gravity:corrections]{第\ref*{theory:gravity:corrections}節}, \hyperref[theory:gravity:geoid]{第\ref*{theory:gravity:geoid}節}, \hyperref[theory:geomagnetism:magnetization]{第\ref*{theory:geomagnetism:magnetization}節}, \hyperref[theory:exploration:observables]{第\ref*{theory:exploration:observables}節}

\par 参考：\cite{gsi_gravity_definitions,kanaya_okuma2011,sato_lei2009_ae}。

\paragraph{関連問題} \hyperref[q:r04-reg-applied-35]{R04-REG-APPLIED-35}

\par 難易度 3（校準初稿）、優先度未評価。

\subsection{令和7年度 問33}
\label{q:r07-reg-applied-33}

\noindent\textbf{原典：} sources/R07/17\_applied\_science.pdf、PDF物理12頁。

\textbf{公式正答：} 5（single）。

\paragraph{問題の要旨} 正規重力と各補正を区別し、密度差によるブーゲー異常の符号を誤る記述を選ぶ。

\Needspace{4\baselineskip}\paragraph{確認した条件・記述（独自要約）}\mbox{}\par

\noindent\textbf{選択方向：} 誤り又は最も不適切・関係の薄いものを1つ\par

\paragraph{問われる核心} 基準密度と重力異常の符号

\paragraph{必要知識} 正規重力、フリーエア・ブーゲー・エトベス補正、密度差の引力積分。

\paragraph{考え方・正答を選ぶ理由（査読待ち初稿）} 測定重力から正規場・高さ・地表質量等の効果を取り除き、地下の密度差を評価する。同じ幾何条件なら重力の差は密度差\(\Delta\rho\)に線形であり、高密度は正の、低密度は負の異常を生む。第5肢はこの符号を逆に記述する。

\Needspace{4\baselineskip}\paragraph{選択肢の検討（査読待ち初稿）}\mbox{}\par

\noindent\textbf{1：適切} 正規重力は地球の質量・自転・形状から導く理論値。基準回転楕円体の外部重力ポテンシャルと自転による遠心ポテンシャルから定める。\par

\noindent\textbf{2：適切} ブーゲー補正は測点と海水準面の間の地殻物質の効果を除く。標高差だけでなく、その間の岩石質量の引力を平板等で近似して除く。\par

\noindent\textbf{3：適切} 船・航空機の移動は遠心効果を変える。地球自転と観測機の相対運動により見掛加速度が変わるためエトベス補正が必要である。\par

\noindent\textbf{4：適切} フリーエア補正は高さの違いの影響を補正する。高さによる重力の鉛直勾配の影響をそろえる操作である。\par

\noindent\textbf{5：不適切} 高密度地域のブーゲー異常は小さく、低密度地域では大きくなる。形状等を同条件にした基準密度との差では、密度の過剰は正、欠損は負の重力異常を生む。第5肢は基本的対応を逆にしている。\par

\paragraph{再利用する考え方} 基準密度と重力異常の符号を、測定量・推定対象・基準面・近似条件に分けて理解する。名称だけで探査法を対応させない。

\paragraph{共通理論} \hyperref[theory:gravity:slab]{第\ref*{theory:gravity:slab}節}, \hyperref[theory:gravity:corrections]{第\ref*{theory:gravity:corrections}節}

\par 参考：\cite{gsi_gravity_definitions}。

\paragraph{関連問題} \hyperref[q:h27-reg-applied-22]{H27-REG-APPLIED-22}

\par 難易度 3（校準初稿）、優先度未評価。

\section{電気・電磁探査}

\subsection{平成23年度 問30}
\label{q:h23-reg-applied-30}

\noindent\textbf{原典：} sources/H23/17\_applied\_science.pdf、PDF物理10頁。

\textbf{公式正答：} 2（single）。

\par\textit{原典本文の校正と独自解説は未作成。分類のみの登録項目。}

\subsection{平成24年度 問20}
\label{q:h24-reg-applied-20}

\noindent\textbf{原典：} sources/H24/17\_applied\_science.pdf、PDF物理7頁。

\textbf{公式正答：} 5（single）。

\par\textit{原典本文の校正と独自解説は未作成。分類のみの登録項目。}

\subsection{平成25年度 問34}
\label{q:h25-reg-applied-34}

\noindent\textbf{原典：} sources/H25/17\_applied\_science.pdf、PDF物理12頁。

\textbf{公式正答：} 3（single）。

\par\textit{原典本文の校正と独自解説は未作成。分類のみの登録項目。}

\subsection{平成25年度 問35}
\label{q:h25-reg-applied-35}

\noindent\textbf{原典：} sources/H25/17\_applied\_science.pdf、PDF物理13頁。

\textbf{公式正答：} 2（single）。

\par\textit{原典本文の校正と独自解説は未作成。分類のみの登録項目。}

\subsection{平成26年度 問22}
\label{q:h26-reg-applied-22}

\noindent\textbf{原典：} sources/H26/17\_applied\_science.pdf、PDF物理9頁。

\textbf{公式正答：} 3（single）。

\par\textit{原典本文の校正と独自解説は未作成。分類のみの登録項目。}

\subsection{平成27年度 問19}
\label{q:h27-reg-applied-19}

\noindent\textbf{原典：} sources/H27/17\_applied\_science.pdf、PDF物理7頁。

\textbf{公式正答：} 3（single）。

\paragraph{問題の要旨} 放射能・地温・電磁・レーダ・重力探査の推定対象を比べ、CSAMTの物性を誤る記述を選ぶ。

\Needspace{4\baselineskip}\paragraph{確認した条件・記述（独自要約）}\mbox{}\par

\noindent\textbf{選択方向：} 誤り又は最も不適切・関係の薄いものを1つ\par

\paragraph{問われる核心} 電磁応答と比抵抗の対応

\paragraph{必要知識} 比抵抗と導電率、平面波インピーダンス、CSAMT遠方場、ガンマ線・温度・重力の観測。

\paragraph{考え方・正答を選ぶ理由（査読待ち初稿）} CSAMTは人工源の遠方場を用いる電磁探査であり、電磁応答から比抵抗構造を推定する。均質・平面波近似なら見掛比抵抗は\(\rho_a=|E/H|^2/(\omega\mu)\)である。磁化率を主対象とする第3肢は物性の対応を誤っている。

\Needspace{4\baselineskip}\paragraph{選択肢の検討（査読待ち初稿）}\mbox{}\par

\noindent\textbf{1：適切} 放射能探査は主にガンマ線強度から断層等の地質情報を得る。天然放射性核種の分布や移動の差を地質情報として利用する。ガンマ線の減衰により観測は主に表層からの寄与であり、深部核種をそのまま透視する意味ではない。\par

\noindent\textbf{2：適切} 地温探査は温泉・地熱・地下水等の調査に用いる。温度分布と熱の移流・伝導を手掛かりに熱水・地下水等を探る。\par

\noindent\textbf{3：不適切} CSAMTは遠方人工源の直交水平電場・磁場から磁化率を得る。主要な推定対象は比抵抗である。水平電場と磁場の比を磁化率そのものと取り違えている。\par

\noindent\textbf{4：適切} レーダは誘電率や導電率の違う境界の電磁反射を捉える。電気特性の差が波の伝播と反射を変えることを利用する。\par

\noindent\textbf{5：適切} 重力は密度構造を求め、標高精度が重要である。標高1 mの誤差は標準鉛直勾配で約0.3086 mGalに相当し、微小重力異常の解釈を妨げる。\par

\paragraph{再利用する考え方} 電磁応答と比抵抗の対応を、測定量・推定対象・基準面・近似条件に分けて理解する。名称だけで探査法を対応させない。

\paragraph{共通理論} \hyperref[theory:exploration:electromagnetic]{第\ref*{theory:exploration:electromagnetic}節}, \hyperref[theory:exploration:observables]{第\ref*{theory:exploration:observables}節}

\par 参考：\cite{ochi_csamt_saijo,gsi_gravity_definitions}。

\par 難易度 3（校準初稿）、優先度未評価。

\subsection{平成28年度 問19}
\label{q:h28-reg-applied-19}

\noindent\textbf{原典：} sources/H28/17\_applied\_science.pdf、PDF物理7頁。

\textbf{公式正答：} 1（single）。

\par\textit{原典本文の校正と独自解説は未作成。分類のみの登録項目。}

\subsection{平成29年度 問32}
\label{q:h29-reg-applied-32}

\noindent\textbf{原典：} sources/H29/17\_applied\_science.pdf、PDF物理11頁。

\textbf{公式正答：} 4（single）。

\par\textit{原典本文の校正と独自解説は未作成。分類のみの登録項目。}

\subsection{令和3年度 問32}
\label{q:r03-reg-applied-32}

\noindent\textbf{原典：} sources/R03/17\_applied\_science.pdf、PDF物理12頁。

\textbf{公式正答：} 4（single）。

\par\textit{原典本文の校正と独自解説は未作成。分類のみの登録項目。}

\section{ボーリング・検層}

\subsection{平成23年度 問31}
\label{q:h23-reg-applied-31}

\noindent\textbf{原典：} sources/H23/17\_applied\_science.pdf、PDF物理10頁。

\textbf{公式正答：} 2（single）。

\par\textit{原典本文の校正と独自解説は未作成。分類のみの登録項目。}

\subsection{平成24年度 問22}
\label{q:h24-reg-applied-22}

\noindent\textbf{原典：} sources/H24/17\_applied\_science.pdf、PDF物理7頁。

\textbf{公式正答：} 2（single）。

\par\textit{原典本文の校正と独自解説は未作成。分類のみの登録項目。}

\subsection{平成27年度 問18}
\label{q:h27-reg-applied-18}

\noindent\textbf{原典：} sources/H27/17\_applied\_science.pdf、PDF物理6頁。

\textbf{公式正答：} 4（single）。

\par\textit{原典本文の校正と独自解説は未作成。分類のみの登録項目。}

\subsection{令和元年度 問33}
\label{q:r01-reg-applied-33}

\noindent\textbf{原典：} sources/R01/17\_applied\_science.pdf、PDF物理12頁。

\textbf{公式正答：} 3（single）。

\par\textit{原典本文の校正と独自解説は未作成。分類のみの登録項目。}

\subsection{令和元年度再試験 問27}
\label{q:r01-ret-applied-27}

\noindent\textbf{原典：} sources/R01\_retest/17\_applied\_science.pdf、PDF物理9頁。

\textbf{公式正答：} 3（single）。

\par\textit{原典本文の校正と独自解説は未作成。分類のみの登録項目。}

\subsection{令和5年度 問34}
\label{q:r05-reg-applied-34}

\noindent\textbf{原典：} sources/R05/17\_applied\_science.pdf、PDF物理12頁。

\textbf{公式正答：} 4（single）。

\par\textit{原典本文の校正と独自解説は未作成。分類のみの登録項目。}

\subsection{令和7年度 問34}
\label{q:r07-reg-applied-34}

\noindent\textbf{原典：} sources/R07/17\_applied\_science.pdf、PDF物理13頁。

\textbf{公式正答：} 4（single）。

\par\textit{原典本文の校正と独自解説は未作成。分類のみの登録項目。}

\section{表面波・微動探査}

\subsection{令和4年度 問34}
\label{q:r04-reg-applied-34}

\noindent\textbf{原典：} sources/R04/17\_applied\_science.pdf、PDF物理12頁。

\textbf{公式正答：} 1（single）。

\par\textit{原典本文の校正と独自解説は未作成。分類のみの登録項目。}

\subsection{令和6年度 問33}
\label{q:r06-reg-applied-33}

\noindent\textbf{原典：} sources/R06/17\_applied\_science.pdf、PDF物理11頁。

\textbf{公式正答：} 2（single）。

\par\textit{原典本文の校正と独自解説は未作成。分類のみの登録項目。}

\subsection{令和7年度 問35}
\label{q:r07-reg-applied-35}

\noindent\textbf{原典：} sources/R07/17\_applied\_science.pdf、PDF物理13頁。

\textbf{公式正答：} 2（single）。

\par\textit{原典本文の校正と独自解説は未作成。分類のみの登録項目。}

\section{地中レーダ}

\subsection{平成24年度 問21}
\label{q:h24-reg-applied-21}

\noindent\textbf{原典：} sources/H24/17\_applied\_science.pdf、PDF物理7頁。

\textbf{公式正答：} 3（single）。

\par\textit{原典本文の校正と独自解説は未作成。分類のみの登録項目。}

\subsection{令和2年度 問33}
\label{q:r02-reg-applied-33}

\noindent\textbf{原典：} sources/R02/17\_applied\_science.pdf、PDF物理12頁。

\textbf{公式正答：} 3（single）。

\par\textit{原典本文の校正と独自解説は未作成。分類のみの登録項目。}

\subsection{令和3年度 問35}
\label{q:r03-reg-applied-35}

\noindent\textbf{原典：} sources/R03/17\_applied\_science.pdf、PDF物理13頁。

\textbf{公式正答：} 4（single）。

\par\textit{原典本文の校正と独自解説は未作成。分類のみの登録項目。}

\subsection{令和5年度 問33}
\label{q:r05-reg-applied-33}

\noindent\textbf{原典：} sources/R05/17\_applied\_science.pdf、PDF物理12頁。

\textbf{公式正答：} 3（single）。

\par\textit{原典本文の校正と独自解説は未作成。分類のみの登録項目。}

\subsection{令和6年度 問35}
\label{q:r06-reg-applied-35}

\noindent\textbf{原典：} sources/R06/17\_applied\_science.pdf、PDF物理12頁。

\textbf{公式正答：} 2（single）。

\par\textit{原典本文の校正と独自解説は未作成。分類のみの登録項目。}

\section{探査信号のフィルタ選択}

\subsection{平成23年度 問28}
\label{q:h23-reg-applied-28}

\noindent\textbf{原典：} sources/H23/17\_applied\_science.pdf、PDF物理9頁。

\textbf{公式正答：} 4（single）。

\par\textbf{検討未確定：} 移動平均の除去は低減であり完全除去ではない。線形性だけで単一応答の畳み込みは保証しない。標本化の原文は一周期と間隔を混ぜた表現であるため、帯域と最高周波数を明示して条件付きで読む。

\paragraph{問題の要旨} 空間フィルタ・平均・インパルス応答・低域通過・標本化を区別する。

\Needspace{4\baselineskip}\paragraph{確認した条件・記述（独自要約）}\mbox{}\par

\noindent\textbf{原典照合：} 否定の対象、指数、周波数と周期、五つの候補の列順を原典画像で照合。\par

\noindent\textbf{前提：} 通常の二次元画像と時系列を対象とする。低域は低周波数すなわち長周期である。\par

\paragraph{問われる核心} 空間フィルタ・平均・インパルス応答・低域通過・標本化を区別する。

\paragraph{必要知識} 通常の二次元画像と時系列を対象とする。低域は低周波数すなわち長周期である。

\paragraph{考え方・正答を選ぶ理由（査読待ち初稿）} 4は低域通過の役割を逆にしている。長周期成分を残し、短周期成分を減衰させるのが低域通過である。他の説明は典型的処理として読むが、完全除去と畳み込みの省略条件を残す。

\Needspace{4\baselineskip}\paragraph{選択肢の検討（査読待ち初稿）}\mbox{}\par

\noindent\textbf{1：適切} 空間フィルタを二次元画像のエッジ・平滑化・差分に用いる。二次元データに対し空間方向の重み付けや差分を行う。一次元・三次元にも拡張でき、二次元しか扱えないという意味ではない。\par

\noindent\textbf{2：条件付きで適切} 注目点付近の移動平均でランダム成分を除去する。独立で平均ゼロの雑音をM点平均すれば分散は1/Mに減る。有限Mで雑音が厳密にゼロになる保証はない。中央窓以外に因果的な後方窓もある。\par

\noindent\textbf{3：条件付きで適切} 線形システムのインパルス応答を重ねて任意入力への応答を求める。線形性により各時刻のインパルスへの応答を重ねることはできる。単一の応答を時間移動する畳み込み表現には時不変性も必要であり、時変系には二時刻の応答を用いる。\par

\noindent\textbf{4：不適切} 低域通過は短周期を残し長周期を減衰させる。周期Tと周波数fは逆数の関係である。低周波数は長周期なので、記述の通過側と減衰側は逆である。\par

\noindent\textbf{5：条件付きで適切} 一周期二標本未満では元データを再構成できない。一般の帯域制限信号を一意に再構成するにはスペクトルの重なりを避ける標本化が必要である。特別な事前情報を持つ信号まで一律に不可能とする定理ではなく、最高周波数と標本化間隔を明示して読む。\par

\paragraph{再利用する考え方} フィルタの名称を周波数で定義してから周期へ変換する。

\paragraph{共通理論} \hyperref[theory:exploration:filters]{第\ref*{theory:exploration:filters}節}, \hyperref[theory:exploration:lti]{第\ref*{theory:exploration:lti}節}, \hyperref[theory:exploration:sampling]{第\ref*{theory:exploration:sampling}節}

\par 参考：\cite{mit_signals2010,mit_moving2012,scipy_notch,scipy_fft,usgs_processing2000}。

\paragraph{関連問題} \hyperref[q:h28-reg-applied-20]{H28-REG-APPLIED-20}

\par 難易度 3（校準初稿）、優先度未評価。

\subsection{平成24年度 問19}
\label{q:h24-reg-applied-19}

\noindent\textbf{原典：} sources/H24/17\_applied\_science.pdf、PDF物理7頁。

\textbf{公式正答：} 5（single）。

\par\textbf{検討未確定：} 移動平均でランダム成分を除去するとは統計的な低減の意味である。信号まで平均化で弱める場合と、相関雑音が残る場合を区別する。

\paragraph{問題の要旨} 相互相関・重合・移動平均・FFT・ノッチの目的を対応させる。

\Needspace{4\baselineskip}\paragraph{確認した条件・記述（独自要約）}\mbox{}\par

\noindent\textbf{原典照合：} 否定の対象、指数、周波数と周期、五つの候補の列順を原典画像で照合。\par

\noindent\textbf{前提：} 二つの時系列の相関、複数記録の平均化、DFT、狭帯域阻止を区別する。\par

\paragraph{問われる核心} 相互相関・重合・移動平均・FFT・ノッチの目的を対応させる。

\paragraph{必要知識} 二つの時系列の相関、複数記録の平均化、DFT、狭帯域阻止を区別する。

\paragraph{考え方・正答を選ぶ理由（査読待ち初稿）} 5が不適切である。ノッチは特定周波数の周囲を阻止するものであり、その周波数だけを通過させる説明は逆である。ランダム雑音の抑制は統計的効果であり有限平均での完全消失ではない。

\Needspace{4\baselineskip}\paragraph{選択肢の検討（査読待ち初稿）}\mbox{}\par

\noindent\textbf{1：適切} 相互相関から二波形の類似と時間遅れを検出する。一方の波形をずらして積の和又は平均を求め、相関の極大から遅れを推定する。極大が常に唯一とは限らず、周期性・雑音に注意する。\par

\noindent\textbf{2：条件付きで適切} 複数記録の重合でランダム雑音を抑える。信号の位相をそろえ、雑音が独立又は無相関なら振幅S/Nは回数の平方根だけ改善する。相関雑音や未補正の走時差がある場合は同じ効果を保証しない。\par

\noindent\textbf{3：条件付きで適切} 注目点付近の移動平均でランダム成分を除去する。独立で平均ゼロの雑音をM点平均すれば分散は1/Mに減る。有限Mで雑音が厳密にゼロになる保証はない。中央窓以外に因果的な後方窓もある。\par

\noindent\textbf{4：適切} FFTで離散時系列のフーリエ係数を高速に計算する。FFTはDFTと異なる変換ではなく、DFTを効率よく計算する算法である。標本化によるaliasingを自動的に取り除く算法ではない。\par

\noindent\textbf{5：不適切} ノッチで指定周波数成分だけを通過させる。ノッチは狭帯域阻止である。通過と遮断を逆にしており、帯域通過の説明と取り違えている。\par

\paragraph{再利用する考え方} 名称ではなく保持する成分と抑える成分を書く。

\paragraph{共通理論} \hyperref[theory:exploration:filters]{第\ref*{theory:exploration:filters}節}, \hyperref[theory:exploration:correlation]{第\ref*{theory:exploration:correlation}節}, \hyperref[theory:exploration:stack]{第\ref*{theory:exploration:stack}節}, \hyperref[theory:exploration:fft]{第\ref*{theory:exploration:fft}節}

\par 参考：\cite{mit_signals2010,mit_moving2012,scipy_notch,scipy_fft,usgs_processing2000}。

\paragraph{関連問題} \hyperref[q:h29-reg-applied-35]{H29-REG-APPLIED-35}

\par 難易度 3（校準初稿）、優先度未評価。

\subsection{平成29年度 問35}
\label{q:h29-reg-applied-35}

\noindent\textbf{原典：} sources/H29/17\_applied\_science.pdf、PDF物理12頁。

\textbf{公式正答：} 5（single）。

\par\textbf{検討未確定：} 移動平均でランダム成分を除去するとは統計的な低減の意味である。信号まで平均化で弱める場合と、相関雑音が残る場合を区別する。

\paragraph{問題の要旨} 相互相関・重合・移動平均・FFT・ノッチの目的を対応させる。

\Needspace{4\baselineskip}\paragraph{確認した条件・記述（独自要約）}\mbox{}\par

\noindent\textbf{原典照合：} 否定の対象、指数、周波数と周期、五つの候補の列順を原典画像で照合。\par

\noindent\textbf{前提：} 二つの時系列の相関、複数記録の平均化、DFT、狭帯域阻止を区別する。\par

\paragraph{問われる核心} 相互相関・重合・移動平均・FFT・ノッチの目的を対応させる。

\paragraph{必要知識} 二つの時系列の相関、複数記録の平均化、DFT、狭帯域阻止を区別する。

\paragraph{考え方・正答を選ぶ理由（査読待ち初稿）} 5が不適切である。ノッチは特定周波数の周囲を阻止するものであり、その周波数だけを通過させる説明は逆である。ランダム雑音の抑制は統計的効果であり有限平均での完全消失ではない。

\Needspace{4\baselineskip}\paragraph{選択肢の検討（査読待ち初稿）}\mbox{}\par

\noindent\textbf{1：適切} 相互相関から二波形の類似と時間遅れを検出する。一方の波形をずらして積の和又は平均を求め、相関の極大から遅れを推定する。極大が常に唯一とは限らず、周期性・雑音に注意する。\par

\noindent\textbf{2：条件付きで適切} 複数記録の重合でランダム雑音を抑える。信号の位相をそろえ、雑音が独立又は無相関なら振幅S/Nは回数の平方根だけ改善する。相関雑音や未補正の走時差がある場合は同じ効果を保証しない。\par

\noindent\textbf{3：条件付きで適切} 注目点付近の移動平均でランダム成分を除去する。独立で平均ゼロの雑音をM点平均すれば分散は1/Mに減る。有限Mで雑音が厳密にゼロになる保証はない。中央窓以外に因果的な後方窓もある。\par

\noindent\textbf{4：適切} FFTで離散時系列のフーリエ係数を高速に計算する。FFTはDFTと異なる変換ではなく、DFTを効率よく計算する算法である。標本化によるaliasingを自動的に取り除く算法ではない。\par

\noindent\textbf{5：不適切} ノッチで指定周波数成分だけを通過させる。ノッチは狭帯域阻止である。通過と遮断を逆にしており、帯域通過の説明と取り違えている。\par

\paragraph{再利用する考え方} 名称ではなく保持する成分と抑える成分を書く。

\paragraph{共通理論} \hyperref[theory:exploration:filters]{第\ref*{theory:exploration:filters}節}, \hyperref[theory:exploration:correlation]{第\ref*{theory:exploration:correlation}節}, \hyperref[theory:exploration:stack]{第\ref*{theory:exploration:stack}節}, \hyperref[theory:exploration:fft]{第\ref*{theory:exploration:fft}節}

\par 参考：\cite{mit_signals2010,mit_moving2012,scipy_notch,scipy_fft,usgs_processing2000}。

\paragraph{関連問題} \hyperref[q:h24-reg-applied-19]{H24-REG-APPLIED-19}

\par 難易度 3（校準初稿）、優先度未評価。

\section{探査の順問題・逆問題と対象雑音}

\subsection{令和7年度 問31}
\label{q:r07-reg-applied-31}

\noindent\textbf{原典：} sources/R07/17\_applied\_science.pdf、PDF物理11頁。

\textbf{公式正答：} 1（single）。

\par\textbf{検討未確定：} 標本化の等号端点と帯域制限、理想再構成の条件が省略される。S/Nの平方根は改善倍率であり、元のS/Nが常に1ではない。

\paragraph{問題の要旨} 順問題と逆問題、対象依存の雑音、標本化と重合を区別する。

\Needspace{4\baselineskip}\paragraph{確認した条件・記述（独自要約）}\mbox{}\par

\noindent\textbf{原典照合：} 否定の対象、指数、周波数と周期、五つの候補の列順を原典画像で照合。\par

\noindent\textbf{前提：} 地下モデルからデータを予測するのが順問題、観測データから地下を推定するのが逆問題である。\par

\paragraph{問われる核心} 順問題と逆問題、対象依存の雑音、標本化と重合を区別する。

\paragraph{必要知識} 地下モデルからデータを予測するのが順問題、観測データから地下を推定するのが逆問題である。

\paragraph{考え方・正答を選ぶ理由（査読待ち初稿）} 1が不適切である。地下構造を求める解析を順問題、計測を逆問題とする対応が逆である。計測そのものはデータ取得であり、順問題とはモデルからデータを計算する写像を指す。

\Needspace{4\baselineskip}\paragraph{選択肢の検討（査読待ち初稿）}\mbox{}\par

\noindent\textbf{1：不適切} 地下推定の解析が順問題、データを得る計測が逆問題である。F(m)でモデルからデータを計算するのが順問題である。与えたデータからmを推定するのが逆問題であり、原典の役割は逆転している。\par

\noindent\textbf{2：適切} 入力信号によって発生した表面波でも解析対象でなければ雑音である。ノイズは対象に対する不要成分の呼び方である。表面波探査では信号となる同じ波が、反射法では相関した妨害波となり得る。\par

\noindent\textbf{3：条件付きで適切} 最高周波数の二倍以上で標本化すれば完全復元できる。帯域制限・理想的標本と再構成を想定した説明である。最高周波数成分を含む場合の等号端点、量子化、有限長、雑音を無視できない。安全な条件はfs>2fmaxであり、境界では位相によって標本が全てゼロとなる例がある。\par

\noindent\textbf{4：適切} ノッチで特定の周波数成分を遮断する。狭い阻止帯域で商用電源などの周期的な混入を抑える。指定帯域だけを通過させるフィルタではない。実装には有限の帯域幅と位相応答がある。\par

\noindent\textbf{5：条件付きで適切} 重合によるランダム雑音に対するS/Nが回数の平方根になる。同一信号の位相と振幅がそろい、雑音が無相関・平均ゼロ・同分散なら成り立つ。一般には元のS/Nに平方根Nを掛ける改善倍率を意味する。\par

\paragraph{再利用する考え方} ノイズをランダムに限らず目的から定義し、写像の向きで順逆を区別する。

\paragraph{共通理論} \hyperref[theory:exploration:inverse]{第\ref*{theory:exploration:inverse}節}, \hyperref[theory:exploration:stack]{第\ref*{theory:exploration:stack}節}, \hyperref[theory:exploration:sampling]{第\ref*{theory:exploration:sampling}節}

\par 参考：\cite{mit_signals2010,mit_moving2012,scipy_notch,scipy_fft,usgs_processing2000,pygimli_inversion}。

\par 難易度 3（校準初稿）、優先度未評価。

\section{逆問題の正則化・最適化と離散化}

\subsection{平成30年度 問35}
\label{q:h30-reg-applied-35}

\noindent\textbf{原典：} sources/H30/17\_applied\_science.pdf、PDF物理12頁。

\textbf{公式正答：} 5（single）。

\par\textbf{検討未確定：} 設問のインバージョンに関係するという広い語には有限要素順モデルも含められる。公式5を算法の直接的役割の比較として保持し、FEMを逆解析から排除しない。

\paragraph{問題の要旨} 逆問題の最適化・正則化と有限要素離散化を区別する。

\Needspace{4\baselineskip}\paragraph{確認した条件・記述（独自要約）}\mbox{}\par

\noindent\textbf{原典照合：} 否定の対象、指数、周波数と周期、五つの候補の列順を原典画像で照合。\par

\noindent\textbf{前提：} モデル更新算法と順モデル計算の離散化を区別する。ただし順モデルは逆問題を構成する内部計算でもある。\par

\paragraph{問われる核心} 逆問題の最適化・正則化と有限要素離散化を区別する。

\paragraph{必要知識} モデル更新算法と順モデル計算の離散化を区別する。ただし順モデルは逆問題を構成する内部計算でもある。

\paragraph{考え方・正答を選ぶ理由（査読待ち初稿）} 公式5は、他四項目が最適化又は逆解の安定化を直接担うのに対し、有限要素法が主に支配方程式の離散化を担うという比較による。ただしFEMで順問題を計算する逆解析は広くある。インバージョンにFEMが関係しないと断定しない。

\Needspace{4\baselineskip}\paragraph{選択肢の検討（査読待ち初稿）}\mbox{}\par

\noindent\textbf{1：適切} マルカート法。非線形最小二乗の更新を減衰項で安定化するLevenberg–Marquardt型の算法である。\par

\noindent\textbf{2：適切} 遺伝的アルゴリズム。モデル候補の集団を評価・選択・変化させて目的関数を探索する最適化法である。大域的最良解を常に保証するわけではない。\par

\noindent\textbf{3：適切} 正則化法。データだけで決まりにくいモデルへ滑らかさや大きさ等の制約を加え、逆問題の不安定性を抑える。\par

\noindent\textbf{4：適切} 共役勾配法。線形化された最小二乗等の連立系又は非線形目的関数の探索に用いる。対称正定値系のCGと非線形CGを区別する。\par

\noindent\textbf{5：条件付きで適切} 有限要素法。直接の最適化・正則化の名称ではなく支配方程式の離散化であるため公式はこれを選ぶ。しかし順モデルや感度計算として逆解析へ用いることができ、無関係という絶対判断は誤りである。\par

\paragraph{再利用する考え方} 目的関数、モデル更新、物理場の数値計算の三つの役割を別々に問う。

\paragraph{共通理論} \hyperref[theory:exploration:inverse]{第\ref*{theory:exploration:inverse}節}, \hyperref[theory:fem:background]{第\ref*{theory:fem:background}節}

\par 参考：\cite{pygimli_inversion,mit_roy_fea1999}。

\paragraph{関連問題} \hyperref[q:r01-reg-applied-35]{R01-REG-APPLIED-35}

\par 難易度 3（校準初稿）、優先度未評価。

\subsection{令和元年度 問35}
\label{q:r01-reg-applied-35}

\noindent\textbf{原典：} sources/R01/17\_applied\_science.pdf、PDF物理12頁。

\textbf{公式正答：} 2（single）。

\paragraph{問題の要旨} フィルタ・差分・共役勾配・平滑化とラテラル検層を区別する。

\Needspace{4\baselineskip}\paragraph{確認した条件・記述（独自要約）}\mbox{}\par

\noindent\textbf{原典照合：} 否定の対象、指数、周波数と周期、五つの候補の列順を原典画像で照合。\par

\noindent\textbf{前提：} 数値処理又はモデル拘束と、孔内で物理量を測る配置を区別する。\par

\paragraph{問われる核心} フィルタ・差分・共役勾配・平滑化とラテラル検層を区別する。

\paragraph{必要知識} 数値処理又はモデル拘束と、孔内で物理量を測る配置を区別する。

\paragraph{考え方・正答を選ぶ理由（査読待ち初稿）} 2が不適切である。ラテラル法は電極配置を用いる電気検層の測定方式である。他四項目は信号処理、方程式の離散化、最適化、解への拘束を担う。

\Needspace{4\baselineskip}\paragraph{選択肢の検討（査読待ち初稿）}\mbox{}\par

\noindent\textbf{1：適切} リニアフィルタ。入力の和と定数倍に対して同じ関係を保つ信号処理であり、畳み込み等の数値計算で実装できる。\par

\noindent\textbf{2：不適切} ラテラル法。USGSの電気検層資料では電極間隔と配置で定義されるlateral logとして確認できる。数値解法の算法を指す語ではない。\par

\noindent\textbf{3：適切} 差分法。微分演算を格子値の差へ置き換えて支配方程式の近似解を求める数値法である。\par

\noindent\textbf{4：適切} 共役勾配法。大型の連立系や最小化を反復して解く数値算法である。\par

\noindent\textbf{5：適切} 平滑化拘束。隣接するモデル値の差などを罰することで、逆解の不要な振動を抑える数値的な条件である。\par

\paragraph{再利用する考え方} 測定の幾何配置とデータ取得後の計算を分ける。

\paragraph{共通理論} \hyperref[theory:exploration:inverse]{第\ref*{theory:exploration:inverse}節}, \hyperref[theory:exploration:seismicterms]{第\ref*{theory:exploration:seismicterms}節}

\par 参考：\cite{usgs_lateral_logging,pygimli_inversion,mit_signals2010}。

\paragraph{関連問題} \hyperref[q:h30-reg-applied-35]{H30-REG-APPLIED-35}

\par 難易度 2（校準初稿）、優先度未評価。

\section{移動平均の雑音低減と周波数零点}

\subsection{平成28年度 問20}
\label{q:h28-reg-applied-20}

\noindent\textbf{原典：} sources/H28/17\_applied\_science.pdf、PDF物理8頁。

\textbf{公式正答：} 2（single）。

\par\textbf{検討未確定：} 2には周期と窓長を合わせれば除去できる具体的反例がある。公式2を保持するが、数学的に周期雑音除去が不可能とは説明しない。3の時不変性と5の帯域・端点条件も省略される。

\paragraph{問題の要旨} 移動平均と周期雑音除去の関係を検討する。

\Needspace{4\baselineskip}\paragraph{確認した条件・記述（独自要約）}\mbox{}\par

\noindent\textbf{原典照合：} 否定の対象、指数、周波数と周期、五つの候補の列順を原典画像で照合。\par

\noindent\textbf{前提：} 通常の移動平均は平滑化を行うが、窓長に応じた周波数応答の零点を持つ。\par

\paragraph{問われる核心} 移動平均と周期雑音除去の関係を検討する。

\paragraph{必要知識} 通常の移動平均は平滑化を行うが、窓長に応じた周波数応答の零点を持つ。

\paragraph{考え方・正答を選ぶ理由（査読待ち初稿）} 公式は2であり、一般的な移動平均を周期雑音専用の選択除去と説明する意図を不適切と読んでいる。しかし周期と窓長が一致する正弦波は移動平均で厳密に消せる。したがって周期雑音を除去できないという一般則へ書き換えず、原典の曖昧さを残す。

\Needspace{4\baselineskip}\paragraph{選択肢の検討（査読待ち初稿）}\mbox{}\par

\noindent\textbf{1：適切} 空間フィルタを二次元画像のエッジ・平滑化・差分に用いる。二次元データに対し空間方向の重み付けや差分を行う。一次元・三次元にも拡張でき、二次元しか扱えないという意味ではない。\par

\noindent\textbf{2：条件付きで適切} 移動平均は時系列の周期的雑音を除去するために用いる。M点等重み平均の周波数応答はsin(Mω/2)/(M sin(ω/2))に位相因子を掛けたものである。ω=2πk/Mの非DC成分は消えるため、窓長を周期へ合わせる反例がある。一般に任意の周期雑音が消えるとは限らない。\par

\noindent\textbf{3：条件付きで適切} 線形システムのインパルス応答を重ねて任意入力への応答を求める。線形性により各時刻のインパルスへの応答を重ねることはできる。単一の応答を時間移動する畳み込み表現には時不変性も必要であり、時変系には二時刻の応答を用いる。\par

\noindent\textbf{4：条件付きで適切} 低域通過で高周波を減衰させ低周波を通す。低域通過の基本的な役割である。無減衰という文言は理想通過帯域のモデルであり、実装の位相・リップルは別に確認する。\par

\noindent\textbf{5：条件付きで適切} 一周期二標本未満では元データを再構成できない。一般の帯域制限信号を一意に再構成するにはスペクトルの重なりを避ける標本化が必要である。特別な事前情報を持つ信号まで一律に不可能とする定理ではなく、最高周波数と標本化間隔を明示して読む。\par

\paragraph{再利用する考え方} 移動平均の効果を目的名で決めず、入力周波数と窓長の応答を計算する。

\paragraph{共通理論} \hyperref[theory:exploration:moving]{第\ref*{theory:exploration:moving}節}, \hyperref[theory:exploration:lti]{第\ref*{theory:exploration:lti}節}, \hyperref[theory:exploration:sampling]{第\ref*{theory:exploration:sampling}節}

\par 参考：\cite{mit_signals2010,mit_moving2012,scipy_notch,scipy_fft,usgs_processing2000}。

\paragraph{関連問題} \hyperref[q:h23-reg-applied-28]{H23-REG-APPLIED-28}

\par 難易度 3（校準初稿）、優先度未評価。

\section{地震探査の走時・画像化用語}

\subsection{平成24年度 問23}
\label{q:h24-reg-applied-23}

\noindent\textbf{原典：} sources/H24/17\_applied\_science.pdf、PDF物理7頁。

\textbf{公式正答：} 4（single）。

\paragraph{問題の要旨} 地震探査の走時・反射用語と重力補正を区別する。

\Needspace{4\baselineskip}\paragraph{確認した条件・記述（独自要約）}\mbox{}\par

\noindent\textbf{原典照合：} 否定の対象、指数、周波数と周期、五つの候補の列順を原典画像で照合。\par

\noindent\textbf{前提：} 地震探査は波の走時・波形、重力探査は加速度の異常を測る。\par

\paragraph{問われる核心} 地震探査の走時・反射用語と重力補正を区別する。

\paragraph{必要知識} 地震探査は波の走時・波形、重力探査は加速度の異常を測る。

\paragraph{考え方・正答を選ぶ理由（査読待ち初稿）} 4のブーゲー補正が不適切である。地震探査に関係する四つは、屈折波の走時切片、初動に現れにくい層、反射を支配するインピーダンス、地形や表層に由来する静的走時ずれに対応する。

\Needspace{4\baselineskip}\paragraph{選択肢の検討（査読待ち初稿）}\mbox{}\par

\noindent\textbf{1：適切} 原点走時。屈折波の走時直線を距離ゼロへ外挿した切片である。直線の傾きと合わせて層厚等を推定する。実測した距離ゼロの初動時刻とは同一視しない。\par

\noindent\textbf{2：適切} ブラインド層。屈折法の初動走時に現れにくい層である。薄い中間層が初動を与えない場合や低速度層を区別し、用語の範囲は資料により異なる。\par

\noindent\textbf{3：適切} 音響インピーダンス。垂直入射の簡単なモデルではZ=ρVであり、境界両側の差が反射振幅を支配する。重力で用いる密度そのものだけとは異なる。\par

\noindent\textbf{4：不適切} ブーゲー補正。測点と基準面の間の質量が重力へ及ぼす効果を補正する重力探査の処理である。地震波走時に対する静補正ではない。\par

\noindent\textbf{5：適切} 静補正。地形や表層低速度層などによる走時ずれを一定の時間シフトとしてそろえる処理である。重力の地形・Bouguer補正と観測量が異なる。\par

\paragraph{再利用する考え方} 処理名だけでなく何を補正し、単位は何かを問う。

\paragraph{共通理論} \hyperref[theory:exploration:seismicterms]{第\ref*{theory:exploration:seismicterms}節}, \hyperref[theory:gravity:corrections]{第\ref*{theory:gravity:corrections}節}

\par 参考：\cite{usgs_refraction2013,epa_seismic_reflection,usgs_static2004,usgs_processing2000}。

\paragraph{関連問題} \hyperref[q:h30-reg-applied-33]{H30-REG-APPLIED-33}

\par 難易度 2（校準初稿）、優先度未評価。

\subsection{平成30年度 問33}
\label{q:h30-reg-applied-33}

\noindent\textbf{原典：} sources/H30/17\_applied\_science.pdf、PDF物理11頁。

\textbf{公式正答：} 5（single）。

\paragraph{問題の要旨} 地震探査の走時・反射用語と重力補正を区別する。

\Needspace{4\baselineskip}\paragraph{確認した条件・記述（独自要約）}\mbox{}\par

\noindent\textbf{原典照合：} 否定の対象、指数、周波数と周期、五つの候補の列順を原典画像で照合。\par

\noindent\textbf{前提：} 地震探査は波の走時・波形、重力探査は加速度の異常を測る。\par

\paragraph{問われる核心} 地震探査の走時・反射用語と重力補正を区別する。

\paragraph{必要知識} 地震探査は波の走時・波形、重力探査は加速度の異常を測る。

\paragraph{考え方・正答を選ぶ理由（査読待ち初稿）} 5のブーゲー補正が不適切である。地震探査に関係する四つは、屈折波の走時切片、初動に現れにくい層、反射を支配するインピーダンス、地形や表層に由来する静的走時ずれに対応する。

\Needspace{4\baselineskip}\paragraph{選択肢の検討（査読待ち初稿）}\mbox{}\par

\noindent\textbf{1：適切} 原点走時。屈折波の走時直線を距離ゼロへ外挿した切片である。直線の傾きと合わせて層厚等を推定する。実測した距離ゼロの初動時刻とは同一視しない。\par

\noindent\textbf{2：適切} ブラインド層。屈折法の初動走時に現れにくい層である。薄い中間層が初動を与えない場合や低速度層を区別し、用語の範囲は資料により異なる。\par

\noindent\textbf{3：適切} 音響インピーダンス。垂直入射の簡単なモデルではZ=ρVであり、境界両側の差が反射振幅を支配する。重力で用いる密度そのものだけとは異なる。\par

\noindent\textbf{4：適切} 静補正。地形や表層低速度層などによる走時ずれを一定の時間シフトとしてそろえる処理である。重力の地形・Bouguer補正と観測量が異なる。\par

\noindent\textbf{5：不適切} ブーゲー補正。測点と基準面の間の質量が重力へ及ぼす効果を補正する重力探査の処理である。地震波走時に対する静補正ではない。\par

\paragraph{再利用する考え方} 処理名だけでなく何を補正し、単位は何かを問う。

\paragraph{共通理論} \hyperref[theory:exploration:seismicterms]{第\ref*{theory:exploration:seismicterms}節}, \hyperref[theory:gravity:corrections]{第\ref*{theory:gravity:corrections}節}

\par 参考：\cite{usgs_refraction2013,epa_seismic_reflection,usgs_static2004,usgs_processing2000}。

\paragraph{関連問題} \hyperref[q:h24-reg-applied-23]{H24-REG-APPLIED-23}

\par 難易度 2（校準初稿）、優先度未評価。

\subsection{令和元年度 問34}
\label{q:r01-reg-applied-34}

\noindent\textbf{原典：} sources/R01/17\_applied\_science.pdf、PDF物理12頁。

\textbf{公式正答：} 4（single）。

\par\textbf{独立解答と公式正答に相違あり：} 独立解答 5。初回独立5対公式4を保持し資料確認後の訂正4を別記する。rigid-boom論文は検索収録の当該原理を確認し、PDF取得は403及びHTML challengeで未達である。日本語表記と装置体系の最終照合を査読対象とする。

\paragraph{問題の要旨} 静補正・深度走時・CMP・リジッドブーム・下方接続を区別する。

\Needspace{4\baselineskip}\paragraph{確認した条件・記述（独自要約）}\mbox{}\par

\noindent\textbf{原典照合：} 否定の対象、指数、周波数と周期、五つの候補の列順を原典画像で照合。\par

\noindent\textbf{前提：} 原典の4はリジットブームシステム、5は下方接続と表記される。下方接続を重力・磁気にしか使わない手法と限定しない。\par

\paragraph{問われる核心} 静補正・深度走時・CMP・リジッドブーム・下方接続を区別する。

\paragraph{必要知識} 原典の4はリジットブームシステム、5は下方接続と表記される。下方接続を重力・磁気にしか使わない手法と限定しない。

\paragraph{考え方・正答を選ぶ理由（査読待ち初稿）} 公式4である。初回5は下方接続をポテンシャル場だけに限定した誤認であった。Stanfordの波動下方接続とmigrationの説明を確認すると、地下へ波動場を外挿する処理も地震探査に関係する。リジッドブームEMIは送受信コイルを用いる電磁誘導計測の配置であり、地震データ処理の他四項目と異なる。

\Needspace{4\baselineskip}\paragraph{選択肢の検討（査読待ち初稿）}\mbox{}\par

\noindent\textbf{1：適切} 静補正。地形や表層低速度層などによる走時ずれを一定の時間シフトとしてそろえる処理である。重力の地形・Bouguer補正と観測量が異なる。\par

\noindent\textbf{2：適切} 深度走時。速度モデルと対応する深度での伝播時間を結ぶ量であり、地震探査の時間・深度対応に用いる。\par

\noindent\textbf{3：適切} CMP重合。共通中点の反射記録を集め、走時差を補正して重ねる反射法の処理である。強い傾斜や不均質では中点が同じ反射点を意味するとは限らない。\par

\noindent\textbf{4：不適切} リジットブームシステム。原典表記を保存する。rigid-boom EMIの一次論文では一次磁場と導体の二次磁場を送受信コイルで扱う電磁計測として記述される。地震波の走時処理ではない。\par

\noindent\textbf{5：適切} 下方接続。地震波動場を深部へ外挿してmigration画像を得る方法に関係する。重力・磁気のポテンシャル場にも下方接続があるため、名前の共通性だけで地震分野から除外しない。\par

\paragraph{再利用する考え方} 同じ数学的処理が複数の物理分野で使われることを許し、計測装置の物理原理を確認する。

\paragraph{共通理論} \hyperref[theory:exploration:seismicterms]{第\ref*{theory:exploration:seismicterms}節}, \hyperref[theory:exploration:continuation]{第\ref*{theory:exploration:continuation}節}, \hyperref[theory:exploration:electromagnetic]{第\ref*{theory:exploration:electromagnetic}節}

\par 参考：\cite{usgs_refraction2013,epa_seismic_reflection,usgs_static2004,usgs_processing2000,stanford_downward2001,rigid_boom2019}。

\paragraph{関連問題} \hyperref[q:r04-reg-applied-32]{R04-REG-APPLIED-32}

\par 難易度 3（校準初稿）、優先度未評価。

\subsection{令和4年度 問32}
\label{q:r04-reg-applied-32}

\noindent\textbf{原典：} sources/R04/17\_applied\_science.pdf、PDF物理11頁。

\textbf{公式正答：} 3（single）。

\paragraph{問題の要旨} 深度走時・migration・静補正・CMPとブーゲー補正を区別する。

\Needspace{4\baselineskip}\paragraph{確認した条件・記述（独自要約）}\mbox{}\par

\noindent\textbf{原典照合：} 否定の対象、指数、周波数と周期、五つの候補の列順を原典画像で照合。\par

\noindent\textbf{前提：} 走時・反射波形の処理と重力の質量補正を分ける。\par

\paragraph{問われる核心} 深度走時・migration・静補正・CMPとブーゲー補正を区別する。

\paragraph{必要知識} 走時・反射波形の処理と重力の質量補正を分ける。

\paragraph{考え方・正答を選ぶ理由（査読待ち初稿）} 3のブーゲー補正が不適切である。migrationは反射・回折記録を地下の反射位置へ戻す処理であり、波動場を下方接続する方法もある。

\Needspace{4\baselineskip}\paragraph{選択肢の検討（査読待ち初稿）}\mbox{}\par

\noindent\textbf{1：適切} 深度走時。速度モデルと深度から伝播時間を対応させる地震探査の量である。\par

\noindent\textbf{2：適切} マイグレーション。傾斜面や回折の見掛け位置を地下の適切な位置へ再配置する画像化処理である。観測時間を一律定数で換算するだけの処理ではない。\par

\noindent\textbf{3：不適切} ブーゲー補正。測点と基準面の間の質量が重力へ及ぼす効果を補正する重力探査の処理である。地震波走時に対する静補正ではない。\par

\noindent\textbf{4：適切} 静補正。地形や表層低速度層などによる走時ずれを一定の時間シフトとしてそろえる処理である。重力の地形・Bouguer補正と観測量が異なる。\par

\noindent\textbf{5：適切} CMP重合。共通中点の記録を集め、NMO等の走時差を整えて重ねる。単に未補正の任意波形を足す方法ではない。\par

\paragraph{再利用する考え方} 同じ補正という語でも観測量と誤差の原因を分ける。

\paragraph{共通理論} \hyperref[theory:exploration:seismicterms]{第\ref*{theory:exploration:seismicterms}節}, \hyperref[theory:exploration:continuation]{第\ref*{theory:exploration:continuation}節}, \hyperref[theory:gravity:corrections]{第\ref*{theory:gravity:corrections}節}

\par 参考：\cite{usgs_refraction2013,epa_seismic_reflection,usgs_static2004,usgs_processing2000,stanford_downward2001}。

\paragraph{関連問題} \hyperref[q:r01-reg-applied-34]{R01-REG-APPLIED-34}

\par 難易度 2（校準初稿）、優先度未評価。

\section{探査の重合とS/N平方根則}

\subsection{平成26年度 問21}
\label{q:h26-reg-applied-21}

\noindent\textbf{原典：} sources/H26/17\_applied\_science.pdf、PDF物理9頁。

\textbf{公式正答：} 3（single）。

\par\textbf{検討未確定：} 平方根則は振幅S/Nと無相関雑音・同位相信号の条件付きであり、群設置も方向・間隔に依存する。

\paragraph{問題の要旨} S/Nの平方根則・群設置・ADC・ノッチ・逆解析を区別する。

\Needspace{4\baselineskip}\paragraph{確認した条件・記述（独自要約）}\mbox{}\par

\noindent\textbf{原典照合：} 否定の対象、指数、周波数と周期、五つの候補の列順を原典画像で照合。\par

\noindent\textbf{前提：} 同じ信号に独立同分散の雑音が重なる場合の振幅S/Nを考える。ADCの標本化周波数は信号帯域に対して設定する。\par

\paragraph{問われる核心} S/Nの平方根則・群設置・ADC・ノッチ・逆解析を区別する。

\paragraph{必要知識} 同じ信号に独立同分散の雑音が重なる場合の振幅S/Nを考える。ADCの標本化周波数は信号帯域に対して設定する。

\paragraph{考え方・正答を選ぶ理由（査読待ち初稿）} 3が不適切である。標本化周波数を最高周波数と同じにすると、一般には異なる連続信号が同じ標本を与える。帯域制限を前提に二倍を超える周波数を用いる。他の候補は代表的な処理を説明する。

\Needspace{4\baselineskip}\paragraph{選択肢の検討（査読待ち初稿）}\mbox{}\par

\noindent\textbf{1：条件付きで適切} 重合のS/N改善が回数の平方根に比例する。同一位相信号はN倍、独立雑音の標準偏差は平方根N倍なので比は平方根N倍となる。\par

\noindent\textbf{2：条件付きで適切} 複数センサーの出力を群として結合し雑音を低減する。独立雑音の平均化に加え、センサーの間隔と入射方向に応じた空間フィルタ効果を利用する。対象の反射波を強めるには位相・配置が適切でなければならない。\par

\noindent\textbf{3：不適切} ADCの標本化周波数を信号の最高周波数と同じにする。一般の信号に必要なNyquist条件を満たさずaliasingを起こす。同じ周波数にするのは正しい一般的設定ではない。\par

\noindent\textbf{4：適切} ノッチで特定の周波数成分を遮断する。狭い阻止帯域で商用電源などの周期的な混入を抑える。指定帯域だけを通過させるフィルタではない。実装には有限の帯域幅と位相応答がある。\par

\noindent\textbf{5：適切} モデルの理論値と観測値の差を非線形最小二乗等で小さくする。順モデルF(m)の残差を目的関数にして地下物性を更新する代表的な逆解析である。線形問題や他のノルムもあり、非線形最小二乗だけが唯一の方法ではない。\par

\paragraph{再利用する考え方} 雑音を振幅で測るか電力で測るかを先に固定する。

\paragraph{共通理論} \hyperref[theory:exploration:stack]{第\ref*{theory:exploration:stack}節}, \hyperref[theory:exploration:sampling]{第\ref*{theory:exploration:sampling}節}, \hyperref[theory:exploration:inverse]{第\ref*{theory:exploration:inverse}節}

\par 参考：\cite{mit_signals2010,mit_moving2012,scipy_notch,scipy_fft,usgs_processing2000,pygimli_inversion}。

\paragraph{関連問題} \hyperref[q:h30-reg-applied-34]{H30-REG-APPLIED-34}

\par 難易度 3（校準初稿）、優先度未評価。

\subsection{平成27年度 問20}
\label{q:h27-reg-applied-20}

\noindent\textbf{原典：} sources/H27/17\_applied\_science.pdf、PDF物理7頁。

\textbf{公式正答：} 1（single）。

\par\textbf{検討未確定：} S/Nの振幅・電力定義が省略される。標本化の等号端点と完全復元、理想帯域通過の実装上の限界を区別する。

\paragraph{問題の要旨} 重合回数・コヒーレンス・標本化・帯域通過・FFTを比較する。

\Needspace{4\baselineskip}\paragraph{確認した条件・記述（独自要約）}\mbox{}\par

\noindent\textbf{原典照合：} 否定の対象、指数、周波数と周期、五つの候補の列順を原典画像で照合。\par

\noindent\textbf{前提：} 振幅S/N、周波数ごとの二信号の関係、理想フィルタを区別する。\par

\paragraph{問われる核心} 重合回数・コヒーレンス・標本化・帯域通過・FFTを比較する。

\paragraph{必要知識} 振幅S/N、周波数ごとの二信号の関係、理想フィルタを区別する。

\paragraph{考え方・正答を選ぶ理由（査読待ち初稿）} 1が不適切である。振幅S/Nは独立雑音のとき平方根N倍であり、N倍ではない。コヒーレンスは相互スペクトルを両信号の自己スペクトルで規格化して比較する。

\Needspace{4\baselineskip}\paragraph{選択肢の検討（査読待ち初稿）}\mbox{}\par

\noindent\textbf{1：不適切} スタックのS/Nが回数Nに比例する。信号はN倍、雑音の振幅尺度は平方根N倍なので振幅S/Nは平方根N倍である。電力比ならN倍となるため比の定義を明示する。\par

\noindent\textbf{2：適切} コヒーレンスは二つの時系列の周波数ごとの類似を示す相関係数である。周波数別の結び付きであり、代表的な二乗コヒーレンスは相互スペクトルの絶対値二乗を自己スペクトルの積で割る。時間領域の遅れ別相関そのものとは区別する。\par

\noindent\textbf{3：条件付きで適切} 最高周波数の二倍以上で標本化すれば完全復元できる。帯域制限・理想的標本と再構成を想定した説明である。最高周波数成分を含む場合の等号端点、量子化、有限長、雑音を無視できない。安全な条件はfs>2fmaxであり、境界では位相によって標本が全てゼロとなる例がある。\par

\noindent\textbf{4：条件付きで適切} 帯域通過は設定帯域を減衰させず、それ以外を遮断する。理想帯域通過の定義として読む。実際の設計では通過帯域のリップル、遷移帯域、有限減衰を持つため完全遮断を一般保証しない。\par

\noindent\textbf{5：適切} FFTは二のべき乗の標本数で効率的である。radix-2では段数がlog2Nとなる。二のべき乗でなければ計算不能という意味ではなく、混合基数やBluestein法もある。\par

\paragraph{再利用する考え方} 振幅比と電力比、理想定義と実装の差を明示する。

\paragraph{共通理論} \hyperref[theory:exploration:stack]{第\ref*{theory:exploration:stack}節}, \hyperref[theory:exploration:correlation]{第\ref*{theory:exploration:correlation}節}, \hyperref[theory:exploration:sampling]{第\ref*{theory:exploration:sampling}節}, \hyperref[theory:exploration:fft]{第\ref*{theory:exploration:fft}節}

\par 参考：\cite{mit_signals2010,mit_moving2012,scipy_notch,scipy_fft,usgs_processing2000,scipy_coherence}。

\par 難易度 3（校準初稿）、優先度未評価。

\subsection{平成30年度 問34}
\label{q:h30-reg-applied-34}

\noindent\textbf{原典：} sources/H30/17\_applied\_science.pdf、PDF物理11頁。

\textbf{公式正答：} 1（single）。

\par\textbf{検討未確定：} 群設置と平方根則は雑音相関・信号整列に依存し、標本化の等号と帯域制限を省略して完全復元を保証できない。

\paragraph{問題の要旨} 重合のS/N指数と電源雑音・群設置・標本化・逆解析を照合する。

\Needspace{4\baselineskip}\paragraph{確認した条件・記述（独自要約）}\mbox{}\par

\noindent\textbf{原典照合：} 否定の対象、指数、周波数と周期、五つの候補の列順を原典画像で照合。\par

\noindent\textbf{前提：} 振幅S/Nは同一位相信号と独立雑音を想定する。\par

\paragraph{問われる核心} 重合のS/N指数と電源雑音・群設置・標本化・逆解析を照合する。

\paragraph{必要知識} 振幅S/Nは同一位相信号と独立雑音を想定する。

\paragraph{考え方・正答を選ぶ理由（査読待ち初稿）} 1が不適切である。平方根Nの向上をNの二乗とする誤りである。電力S/NでもN倍でありNの二乗にはならない。周期雑音には狭帯域阻止、物性推定には順モデルの残差最小化を対応させる。

\Needspace{4\baselineskip}\paragraph{選択肢の検討（査読待ち初稿）}\mbox{}\par

\noindent\textbf{1：不適切} スタックのS/N改善が回数Nの二乗に比例する。信号の振幅Nに対して雑音の標準偏差は平方根Nである。振幅比も電力比もNの二乗とはならない。\par

\noindent\textbf{2：条件付きで適切} 複数センサーの出力を群として結合し雑音を低減する。独立雑音の平均化に加え、センサーの間隔と入射方向に応じた空間フィルタ効果を利用する。対象の反射波を強めるには位相・配置が適切でなければならない。\par

\noindent\textbf{3：適切} ノッチで特定の周波数成分を遮断する。狭い阻止帯域で商用電源などの周期的な混入を抑える。指定帯域だけを通過させるフィルタではない。実装には有限の帯域幅と位相応答がある。\par

\noindent\textbf{4：条件付きで適切} 最高周波数の二倍以上で標本化すれば完全復元できる。帯域制限・理想的標本と再構成を想定した説明である。最高周波数成分を含む場合の等号端点、量子化、有限長、雑音を無視できない。安全な条件はfs>2fmaxであり、境界では位相によって標本が全てゼロとなる例がある。\par

\noindent\textbf{5：適切} モデルの理論値と観測値の差を非線形最小二乗等で小さくする。順モデルF(m)の残差を目的関数にして地下物性を更新する代表的な逆解析である。線形問題や他のノルムもあり、非線形最小二乗だけが唯一の方法ではない。\par

\paragraph{再利用する考え方} 加算した信号と雑音を別々に計算して指数の誤りを検出する。

\paragraph{共通理論} \hyperref[theory:exploration:stack]{第\ref*{theory:exploration:stack}節}, \hyperref[theory:exploration:filters]{第\ref*{theory:exploration:filters}節}, \hyperref[theory:exploration:inverse]{第\ref*{theory:exploration:inverse}節}

\par 参考：\cite{mit_signals2010,mit_moving2012,scipy_notch,scipy_fft,usgs_processing2000,pygimli_inversion}。

\paragraph{関連問題} \hyperref[q:h26-reg-applied-21]{H26-REG-APPLIED-21}, \hyperref[q:r04-reg-applied-33]{R04-REG-APPLIED-33}

\par 難易度 3（校準初稿）、優先度未評価。

\subsection{令和4年度 問33}
\label{q:r04-reg-applied-33}

\noindent\textbf{原典：} sources/R04/17\_applied\_science.pdf、PDF物理11頁。

\textbf{公式正答：} 5（single）。

\par\textbf{検討未確定：} 完全復元の等号端点・帯域制限、群配置の方向と間隔を要する。表面波抑制を独立雑音の平方根則だけで説明しない。

\paragraph{問題の要旨} モデル更新・ノッチ・受振器群・標本化・垂直重合を区別する。

\Needspace{4\baselineskip}\paragraph{確認した条件・記述（独自要約）}\mbox{}\par

\noindent\textbf{原典照合：} 否定の対象、指数、周波数と周期、五つの候補の列順を原典画像で照合。\par

\noindent\textbf{前提：} 同一起振点の反復で対象反射の位相をそろえる。群配置は対象波と妨害波の空間応答を考える。\par

\paragraph{問われる核心} モデル更新・ノッチ・受振器群・標本化・垂直重合を区別する。

\paragraph{必要知識} 同一起振点の反復で対象反射の位相をそろえる。群配置は対象波と妨害波の空間応答を考える。

\paragraph{考え方・正答を選ぶ理由（査読待ち初稿）} 5が不適切である。垂直重合の振幅S/Nは理論上平方根N倍でありNの二乗ではない。表面波は反射波に対する妨害波となるが、空間的に相関した波なので独立ランダム雑音の平均化と同じ説明だけでは足りない。

\Needspace{4\baselineskip}\paragraph{選択肢の検討（査読待ち初稿）}\mbox{}\par

\noindent\textbf{1：適切} モデルの理論値と観測値の差を非線形最小二乗等で小さくする。順モデルF(m)の残差を目的関数にして地下物性を更新する代表的な逆解析である。線形問題や他のノルムもあり、非線形最小二乗だけが唯一の方法ではない。\par

\noindent\textbf{2：適切} ノッチで特定の周波数成分を遮断する。狭い阻止帯域で商用電源などの周期的な混入を抑える。指定帯域だけを通過させるフィルタではない。実装には有限の帯域幅と位相応答がある。\par

\noindent\textbf{3：条件付きで適切} 受振器群で表面波を抑え反射波を強めランダム雑音を減らす。複数受振器の結合には空間応答があり、反射と表面波の見掛け速度・波長に応じた選別を行える。配列が悪ければ必要信号も弱まるため無条件ではない。\par

\noindent\textbf{4：条件付きで適切} 最高周波数の二倍以上で標本化すれば完全復元できる。帯域制限・理想的標本と再構成を想定した説明である。最高周波数成分を含む場合の等号端点、量子化、有限長、雑音を無視できない。安全な条件はfs>2fmaxであり、境界では位相によって標本が全てゼロとなる例がある。\par

\noindent\textbf{5：不適切} 垂直スタックのS/Nが回数Nの二乗に比例する。独立同分散雑音の振幅S/N改善は平方根N倍である。Nの二乗は信号電力だけの増加を雑音の増加と比較せず取り出した誤りである。\par

\paragraph{再利用する考え方} 信号の種類、統計的相関、配列の空間応答を分ける。

\paragraph{共通理論} \hyperref[theory:exploration:stack]{第\ref*{theory:exploration:stack}節}, \hyperref[theory:exploration:filters]{第\ref*{theory:exploration:filters}節}, \hyperref[theory:exploration:inverse]{第\ref*{theory:exploration:inverse}節}

\par 参考：\cite{mit_signals2010,mit_moving2012,scipy_notch,scipy_fft,usgs_processing2000,pygimli_inversion}。

\paragraph{関連問題} \hyperref[q:h30-reg-applied-34]{H30-REG-APPLIED-34}

\par 難易度 3（校準初稿）、優先度未評価。

\section{探査波動の伝播原理}

\subsection{令和2年度 問35}
\label{q:r02-reg-applied-35}

\noindent\textbf{原典：} sources/R02/17\_applied\_science.pdf、PDF物理12頁。

\textbf{公式正答：} 2（single）。

\par\textbf{検討未確定：} 相反性と重ね合わせは媒体・線形性の条件付きである。等価層は異なる分野で使われる語でもあるため、ポテンシャル場の基本的用法と波動伝播の原理を区別する。

\paragraph{問題の要旨} 弾性波の伝播原理とポテンシャル場の等価層を区別する。

\Needspace{4\baselineskip}\paragraph{確認した条件・記述（独自要約）}\mbox{}\par

\noindent\textbf{原典照合：} 否定の対象、指数、周波数と周期、五つの候補の列順を原典画像で照合。\par

\noindent\textbf{前提：} 線形で相反的な媒体における波動の重ね合わせ・波面構成・伝播経路を対象とする。\par

\paragraph{問われる核心} 弾性波の伝播原理とポテンシャル場の等価層を区別する。

\paragraph{必要知識} 線形で相反的な媒体における波動の重ね合わせ・波面構成・伝播経路を対象とする。

\paragraph{考え方・正答を選ぶ理由（査読待ち初稿）} 2の等価層の原理が不適切である。USGSの等価層は重力の密度層や磁気の双極子層でポテンシャル場を表現するものとして確認できる。これは地震波伝播を記述する基本原理の四項目とは役割が異なる。

\Needspace{4\baselineskip}\paragraph{選択肢の検討（査読待ち初稿）}\mbox{}\par

\noindent\textbf{1：適切} フェルマの原理。波線経路に対する走時の停留条件を述べる。あらゆる経路が絶対的最短時間になるという一般化は避ける。\par

\noindent\textbf{2：不適切} 等価層の原理。重力・磁気の場を等価な密度又は双極子の層で表す考え方である。弾性波の波面や伝播経路を定める基本則とは異なる。\par

\noindent\textbf{3：適切} 相反原理。相反性が成立する線形媒体では震源と受振点を交換した応答が対応する。媒体や成分の条件を保つ必要がある。\par

\noindent\textbf{4：適切} ホイヘンスの原理。波面上の点からの二次波の構成で次の波面を考える。地震波の伝播・画像化の理解に関係する。\par

\noindent\textbf{5：条件付きで適切} 重ね合わせの原理。線形の波動方程式では解の線形結合も解となる。破壊を伴う大きな非線形変形へ無条件に適用しない。\par

\paragraph{再利用する考え方} 原理の名前を、場の表現・波面・経路・線形性・相反性という役割へ対応させる。

\paragraph{共通理論} \hyperref[theory:exploration:principles]{第\ref*{theory:exploration:principles}節}

\par 参考：\cite{usgs_equivalent2007,stanford_downward2001,mit_signals2010}。

\paragraph{関連問題} \hyperref[q:r07-reg-applied-32]{R07-REG-APPLIED-32}

\par 難易度 3（校準初稿）、優先度未評価。

\subsection{令和7年度 問32}
\label{q:r07-reg-applied-32}

\noindent\textbf{原典：} sources/R07/17\_applied\_science.pdf、PDF物理12頁。

\textbf{公式正答：} 3（single）。

\par\textbf{検討未確定：} 重ね合わせは線形媒体に限る。スネルの法則は波種と境界条件を指定して用いる。

\paragraph{問題の要旨} ホイヘンス・重ね合わせ・レンツ・スネル・フェルマを区別する。

\Needspace{4\baselineskip}\paragraph{確認した条件・記述（独自要約）}\mbox{}\par

\noindent\textbf{原典照合：} 否定の対象、指数、周波数と周期、五つの候補の列順を原典画像で照合。\par

\noindent\textbf{前提：} 地震探査で扱う弾性波の波面、線形性、屈折、波線の走時を対象とする。\par

\paragraph{問われる核心} ホイヘンス・重ね合わせ・レンツ・スネル・フェルマを区別する。

\paragraph{必要知識} 地震探査で扱う弾性波の波面、線形性、屈折、波線の走時を対象とする。

\paragraph{考え方・正答を選ぶ理由（査読待ち初稿）} 3のレンツの法則は電磁誘導が磁束変化を妨げる向きに生じることを表す。電磁探査には関係するが、他四項目のような弾性波伝播の基本則ではない。

\Needspace{4\baselineskip}\paragraph{選択肢の検討（査読待ち初稿）}\mbox{}\par

\noindent\textbf{1：適切} ホイヘンスの原理。二次波で波面を構成する考え方であり、波動伝播の理解に関係する。\par

\noindent\textbf{2：条件付きで適切} 重ね合わせの原理。線形波動では複数波の和を解として扱うことができる。大変形や非線形な破壊過程は同じ前提ではない。\par

\noindent\textbf{3：不適切} レンツの法則。電磁誘導の向きを磁束変化との関係から決める法則であり、地震波の反射・屈折の基本則とは異なる。\par

\noindent\textbf{4：適切} スネルの法則。境界の位相整合により角度と波速を結ぶ。地震波ではP波・S波へのモード変換を含む各波速で考える。\par

\noindent\textbf{5：適切} フェルマの原理。波線の走時が経路変分に対して停留する条件を用いる。最小という単語だけで複数経路の可能性を捨てない。\par

\paragraph{再利用する考え方} 電磁気の誘導と弾性波の伝播を、扱う物理量から分ける。

\paragraph{共通理論} \hyperref[theory:exploration:principles]{第\ref*{theory:exploration:principles}節}, \hyperref[theory:optics:refraction]{第\ref*{theory:optics:refraction}節}

\par 参考：\cite{mit_signals2010,stanford_downward2001,epa_seismic_reflection}。

\paragraph{関連問題} \hyperref[q:r02-reg-applied-35]{R02-REG-APPLIED-35}

\par 難易度 3（校準初稿）、優先度未評価。

\section{未確定}

\subsection{平成25年度 問33}
\label{q:h25-reg-applied-33}

\noindent\textbf{原典：} sources/H25/17\_applied\_science.pdf、PDF物理12頁。

\textbf{公式正答：} 5（single）。

\par\textit{原典本文の校正と独自解説は未作成。分類のみの登録項目。}

\subsection{平成26年度 問19}
\label{q:h26-reg-applied-19}

\noindent\textbf{原典：} sources/H26/17\_applied\_science.pdf、PDF物理8頁。

\textbf{公式正答：} 5（single）。

\par\textit{原典本文の校正と独自解説は未作成。分類のみの登録項目。}

\subsection{平成26年度 問23}
\label{q:h26-reg-applied-23}

\noindent\textbf{原典：} sources/H26/17\_applied\_science.pdf、PDF物理10頁。

\textbf{公式正答：} 4（single）。

\par\textit{原典本文の校正と独自解説は未作成。分類のみの登録項目。}

\subsection{平成27年度 問21}
\label{q:h27-reg-applied-21}

\noindent\textbf{原典：} sources/H27/17\_applied\_science.pdf、PDF物理8頁。

\textbf{公式正答：} 1（single）。

\par\textit{原典本文の校正と独自解説は未作成。分類のみの登録項目。}

\subsection{平成28年度 問21}
\label{q:h28-reg-applied-21}

\noindent\textbf{原典：} sources/H28/17\_applied\_science.pdf、PDF物理8頁。

\textbf{公式正答：} 1（single）。

\par\textit{原典本文の校正と独自解説は未作成。分類のみの登録項目。}

\subsection{平成28年度 問22}
\label{q:h28-reg-applied-22}

\noindent\textbf{原典：} sources/H28/17\_applied\_science.pdf、PDF物理9頁。

\textbf{公式正答：} 1（single）。

\par\textit{原典本文の校正と独自解説は未作成。分類のみの登録項目。}

\subsection{平成28年度 問23}
\label{q:h28-reg-applied-23}

\noindent\textbf{原典：} sources/H28/17\_applied\_science.pdf、PDF物理9頁。

\textbf{公式正答：} 3（single）。

\par\textit{原典本文の校正と独自解説は未作成。分類のみの登録項目。}

\subsection{平成29年度 問33}
\label{q:h29-reg-applied-33}

\noindent\textbf{原典：} sources/H29/17\_applied\_science.pdf、PDF物理11頁。

\textbf{公式正答：} 3（single）。

\par\textit{原典本文の校正と独自解説は未作成。分類のみの登録項目。}

\subsection{令和元年度 問32}
\label{q:r01-reg-applied-32}

\noindent\textbf{原典：} sources/R01/17\_applied\_science.pdf、PDF物理11頁。

\textbf{公式正答：} 5（single）。

\par\textit{原典本文の校正と独自解説は未作成。分類のみの登録項目。}

\subsection{令和2年度 問34}
\label{q:r02-reg-applied-34}

\noindent\textbf{原典：} sources/R02/17\_applied\_science.pdf、PDF物理12頁。

\textbf{公式正答：} 4（single）。

\par\textit{原典本文の校正と独自解説は未作成。分類のみの登録項目。}

\subsection{令和5年度 問35}
\label{q:r05-reg-applied-35}

\noindent\textbf{原典：} sources/R05/17\_applied\_science.pdf、PDF物理13頁。

\textbf{公式正答：} 2（single）。

\par\textit{原典本文の校正と独自解説は未作成。分類のみの登録項目。}
